\section*{附录习题}
\phantomsection
\addcontentsline{toc}{section}{附录习题}

\begin{enumerate}
\def\labelenumi{\arabic{enumi})}
\tightlist
\item
  设 I 是一个非空的右滤预序集，并设 \((\mathbf{A}_{\alpha}, \varphi_{\beta\alpha})\) 是关于 I 的环的归纳系统。对于 I 中的每个指标 \(\alpha\)，给定 \(q_{\alpha}\)（\(A_{\alpha}\) 的一个理想）。假设 \(\varphi_{\beta\alpha}(q_{\alpha}) A_{\beta} = q_{\beta}\) 对于 \(\beta \geqslant \alpha\) 成立。记 A 为 \(A_{\alpha}\) 的归纳极限，并且对于 \(\alpha \in I\)，记 \(\varphi_{\alpha}: A_{\alpha} \to A\) 为典范同态。置
\end{enumerate}

\[
\mathfrak {q} = \varinjlim \mathfrak {q} _ {\alpha}.
\]

\begin{enumerate}
\def\labelenumi{\alph{enumi})}
\item
  若对于每个 \(\alpha \in I\)，\(q_{\alpha}\) 是 \(A_{\alpha}\) 的极大理想，则 \(q\) 是 \(A\) 的极大理想。
\item
  若对于每个 \(\alpha \in I\)，\(\mathfrak{q}_{\alpha}\) 包含于 \(A_{\alpha}\) 的根基中，则 \(\mathfrak{q}\) 包含于 \(A\) 的根基中。
\item
  若对于每个 \(\alpha\in\mathbf{I}\) 和每个 \(\beta\in\mathbf{I}\)，当 \(\beta\geqslant\alpha\) 时，同态 \(\varphi_{\beta\alpha}\) 都是忠实平坦的，则对于所有 \(\alpha\in\mathbf{I}\)，\(\varphi_{\alpha}\) 都是忠实平坦的。
\end{enumerate}

以下设 c) 的假设成立。

\(d)\) 若对于每个 \(\alpha\in I\)，\(\mathbf{A}_{\alpha}\) 关于 \(\mathfrak{q}_{\alpha}\)-进拓扑是分离的，则 A 关于 \(\mathfrak{q}\)-进拓扑是分离的。

\begin{enumerate}
\def\labelenumi{\alph{enumi})}
\setcounter{enumi}{4}
\tightlist
\item
  若对于每个 \(\alpha \in I\)，\(A_{\alpha}\) 是诺特的，且 A/q 是诺特的，则 \((1 + q)^{-1}A\) 是诺特的。（理由同附录命题 1 的证明，此外使用下面的结果，该结果将被证明：
\end{enumerate}

设 B 是一个交换环，p 是 B 的一个理想，\(\widehat{B}\) 是 B 关于 p-进拓扑的分离完备化，且 \(i: B \to \widehat{B}\) 是典范映射。则存在唯一一个环同态 \(j: (1 + p)^{-1}B \to \widehat{B}\)，它延拓 \(i\)。对于 \((1 + p)^{-1}B\) 的每个极大理想 m，都有 \(j(m)\widehat{B} \neq \widehat{B}\)。

\begin{enumerate}
\def\labelenumi{\arabic{enumi})}
\setcounter{enumi}{1}
\item
  设 A 是一个局部环，B 是 A 的一次爆破。如果 A 关于由极大理想 \(m_{A}\) 定义的拓扑是分离的，则 B 关于由极大理想 \(m_{B}\) 定义的拓扑也是分离的。
\item
  设 \(k\) 是特征 \(p > 0\) 的非完美域，且 \(\mathbf{R} = k[\mathrm{T}]_{(\mathrm{T})}\)。设 \(a\) 是 \(k\) 中不是 \(p\) 次幂的元素。证明环 \(\mathbf{A} = \mathbf{R}[\mathbf{X}] / (\mathbf{X}^p - a\mathbf{T}^p)\) 是局部环、整环且剩余域为 \(k\)。证明环 \(\mathbf{B} = \mathbf{A}[\mathbf{Y}] / (\mathbf{Y}^p - a)\)（它是 \(\mathbf{A}\) 的一次爆破）不是约化的。
\item
  设 A 是赋有赋值 \(v_{\mathbf{A}}\) 的局部环。设 B 是 A 的一次爆破。证明赋值 \(v_{\mathbf{A}}\) 唯一延拓为 B 的赋值 \(v_{\mathbf{B}}\)，且值群 \(v_{\mathbf{A}}(\mathbf{A})\) 和 \(v_{\mathbf{B}}(\mathbf{B})\) 相同。如果 A 是带一致化元 \(\pi\) 的离散赋值环，则 B 也是离散赋值环，其一致化元之一是 \(\pi\) 在 B 中的像。
\item
  设 A 是剩余域为 \(\kappa(A)\) 的局部环。设 B 是 A 的一次爆破，其剩余域为 \(\kappa(B)\)。
\end{enumerate}

\begin{enumerate}
\def\labelenumi{\alph{enumi})}
\item
  若 \(\kappa(\mathbf{B}) = \kappa(\mathbf{A})\)，则 \(\mathbf{A} = \mathbf{B}\)。
\item
  若 \(\kappa(\mathbf{B})\) 是 \(\kappa(\mathbf{A})\) 的代数扩张，则 A-代数 B 是整的。
\item
  若 \(\kappa(B)\) 在 \(\kappa(A)\) 上的次数有限，则 \(B\) 是自由 A-模，其秩等于 \(\kappa(B)\) 在 \(\kappa(A)\) 上的次数。
\end{enumerate}

\begin{enumerate}
\def\labelenumi{\arabic{enumi})}
\setcounter{enumi}{5}
\tightlist
\item
  设 A 是剩余域为 \(k\) 的局部环。设 I 是指标集。证明环 \(\mathbf{A}](\mathbf{X}_i)_{i\in \mathbf{I}}[\)（参见附录，例 2）是 A 的一次爆破，其剩余域是 \(k((\mathbf{X}_i)_{i\in \mathbf{I}})\)（\(k\) 的纯超越扩张）。
\end{enumerate}

